% ============================================================================
% DISCUSSION. Three paragraphs, each with one job:
%   1 synthesis and the actual claim boundary
%   2 limitations of the CONTROL MACHINERY, four kept distinct:
%       support -> approximate h_phi -> surrogate -> inference compute
%   3 limitations of the MOLECULAR REPRESENTATION, then the implication
% Deliberate: para 1 concedes that future-aware control is not unique to
% COMPOSE. The claim is the substrate and the decomposition, not the Doob
% transform. Do not reinstate 'endpoint generation followed by scoring or
% repair' -- that was a false dichotomy. The closing sentence says
% 'interpretable states and explicit executable transformations' rather than
% 'valid states', which would sound universal and theorem-like.
% ============================================================================
\section{Discussion and limitations}
\label{sec:discussion}

Taken together, our results support treating molecular design as control of a learned executable process rather than as a sequence of independent endpoint-generation problems. \method{} separates the support of molecular change from the probability assigned to those changes and from the task-specific objective used to control them. The learned reference process captures state-dependent preferences over executable successors, while future reachability can improve multi-step decisions without modifying that process. Because the states being controlled are themselves complete molecules, the same formulation also permits a realized intermediate to be reused after an objective change and allows statewise requirements to act directly on the trajectory. Future-aware generative control is not unique to \method{}; the distinction is the substrate on which it operates, a learned canonical transition process over executable molecular graphs whose support, plausibility, and task-specific purpose remain separate.

This separation also makes the principal limitations clear. Control can only redistribute probability over futures represented by the executable support and assigned sufficient mass by $R_\theta$; neither exact nor learned control can recover a useful route that the underlying molecular process does not contain. The exact finite-horizon guarantees further rely on exact backward values, whereas molecular-scale control uses an approximate $h_\phi$. Empirical improvements with learned control therefore establish useful future-aware steering, not exact terminal conditioning at scale. Property-directed control also inherits calibration and extrapolation error from the property models used to define desirability, and $R_\theta$ should not be interpreted as physical kinetics or synthetic-route likelihood. Finally, richer control can require substantial inference-time computation through continuation evaluation or particle methods; our benchmark results establish task performance rather than a general claim of computational efficiency.

The declared molecular state space is itself deliberately limited. Connectivity and valence admissibility in a 2D molecular graph do not guarantee stereochemical correctness, three-dimensional plausibility, synthesizability, stability, or biological activity, and some current transition families rely on synthetic rather than data-backed supervision. Extending \method{} therefore requires improving the molecular process as well as the controller: richer executable state spaces, broader data-backed transition supervision, and higher-fidelity physical or experimental feedback could expand which useful molecular futures are both represented and preferred. More broadly, the framework suggests that when a domain admits interpretable states and explicit executable transformations between them, generation can be formulated as learning a reusable stochastic process over those transformations and adapting its control as the purpose changes.
